\documentclass[11pt]{article}
\usepackage[letterpaper,hmargin=1.02in,vmargin=0.93in]{geometry}
\usepackage[T1]{fontenc}
\usepackage{lmodern,amsmath,amssymb,amsthm,microtype}
\usepackage[hidelinks]{hyperref}
\newtheorem{theorem}{Theorem}
\newtheorem{lemma}[theorem]{Lemma}
\newtheorem{proposition}[theorem]{Proposition}
\newcommand{\F}{\mathbb F}
\newcommand{\NL}{\mathsf{NL}}
\newcommand{\US}{\mathsf{USPACE}}
\newcommand{\calC}{\mathcal C}
\newcommand{\calF}{\mathcal F}
\newcommand{\fr}{\operatorname{Fr}}
\newcommand{\id}{\operatorname{id}}
\newcommand{\pr}{\Pr}
\allowdisplaybreaks

\hypersetup{pdftitle={Unambiguous reachability in log n times log log n space}}
\title{Unambiguous reachability in \(O(\log n\log\log n)\) space}
\author{}
\date{November 2026}
\begin{document}
\maketitle
\begin{abstract}
We give an explicit family of additive weights for layered directed graphs.
For every graph on $N$ vertices, at least one assignment gives a unique
minimum path between every reachable pair. The seed length, the bit
length of a path weight, and the space for evaluating the assignment are
all $O(\log N\log\log N)$. Exhaustively searching the family gives an
unambiguous algorithm for directed reachability in the same space.
\end{abstract}

\section{Statement and proof strategy}

\begin{theorem}\label{thm:main}
$\NL\subseteq\US(O(\log n\log\log n))$.
\end{theorem}
We give the proof for a simple layered DAG on $N$ vertices. Layers are numbered
$0,\ldots,d$ with $d\le N$, every edge advances one layer, and vertices have
distinct identifiers in $[N]$. The standard logspace reduction from $\NL$
gives graphs of this form with $N$ polynomial in the input length. Write
$L=\lceil\log_2 N\rceil$ and $\ell=\lceil\log_2(2L)\rceil$. All bounds below
are for sufficiently large $N$. If the graph contains no path of length two,
reachability can be decided directly in logarithmic space by testing equality
and adjacency. We henceforth assume at least one such path exists.

An assignment is \emph{min-isolating} if every reachable pair has a unique
minimum-weight path. We minimize a sequence of additive vertex weights
lexicographically. Optimal subpaths are optimal for their own endpoints,
by replacement. When comparing paths between fixed endpoints, we omit the
endpoints' contributions. The entire vector objective can be encoded by
one integer weight.


We use the isolation framework of Reinhardt and Allender~\cite{RA} and
the shifted hashes of van Melkebeek and Prakriya~\cite{vMP}. The binary
layer decomposition also appears in Hoza~\cite[Appendix A]{Hoza}.
The seed in our construction can depend on the graph. Our running time
includes an exhaustive seed search.

The proof keeps track of the vertices of a path on the layers whose
binary valuations lie in the highest occupied band of a window.
We call this ordered list its skeleton. Doubling the width of the band
exposes an old skeleton, followed by old skeletons in the gaps. An
entire traversal of these exposures corresponds to a single path, even
when alternative answers to individual questions come from different paths.

There are two uses of this distinction. First, we bound the number of lists of
subproblems that a traversal can produce. In a fixed plan for naming
their endpoints, each path makes only $O(L)$ midpoint choices, and any
physical layer occurs at most $\ell$ times. If many choices survive
a perturbation, their alternatives certify many independent linear
equations. Repeating the perturbation reduces the logarithm of the
number of feasible transcripts from $O(L\ell)$ to $O(L)$.

Second, we show that the number of skeletons on a whole list of
subproblems remains small as the band width doubles. The separate
perturbations at these scales share bit positions in a single additive
weight. At each step we keep every path that could minimize the cost
under a bounded future perturbation. This makes the surviving classes
nested while retaining the linear equations needed to bound ambiguity.
\begin{lemma}[Small-bias labels]\label{lem:bias}
For a prime $p$, $n_0$ coordinates, and $0<\varepsilon<1$, there is an explicit
distribution on $h\in\F_p^{n_0}$ with every nontrivial additive character
of expectation of absolute value at most $\varepsilon$. A seed and the
space to evaluate a coordinate each use $O(\log n_0+\log p+\log(1/\varepsilon))$
bits. For any fixed system of $m$ linearly independent affine constraints
on $h$, the probability of satisfying all of them is at most
$p^{-m}+\varepsilon$.
\end{lemma}
\begin{proof}
Choose the smallest extension $E=\F_{p^u}$ with $|E|>n_0/\varepsilon$.
For uniform independent $x,y\in E$ and any nonzero
$\F_p$-linear functional $\pi:E\to\F_p$, put $h_j=\pi(yx^j)$, $1\le j\le n_0$.
A nontrivial linear combination has a nonzero coefficient polynomial
of degree at most $n_0$ in $x$. Away from its roots the linear combination
is uniform by the choice of $y$. This proves the bias bound.
Expanding the indicator of the affine system in characters gives a constant term
$p^{-m}$, and the remaining coefficients have total absolute value at most one.
For uniform evaluation, enumerate monic degree-$u$ polynomials and test
irreducibility by trial division. In the first resulting field representation
one may take $\pi$ to be the constant coefficient. The search and the field
arithmetic use $O(u\log p)$ space. Here seed length means the length
of the encoding of a point of the support. For odd $p$ we use uniform
field elements in the existence proof and skip invalid encodings in
the exhaustive search.
\end{proof}

\section{Thinning the feasible midpoints}

For two vertices $x,y$ in layers $a<b$ with $b-a\ge 2$, let their
\emph{pivot} be the internal layer $i$ maximizing $\nu_2(i)$.
The pivot is unique because between two positive integers of the same
$2$-adic valuation lies an integer of higher valuation. The feasible
midpoints for $(x,y)$ are the vertices in its pivot layer on an
$x$--$y$ path.

\begin{lemma}\label{lem:raw}
There is an explicitly enumerable family of nonnegative integer vertex
weights $V$ for $N$-vertex layered graphs with the following property.
For each graph, some $V$ in the family ensures that, for every reachable
pair $(x,y)$ having a pivot, at most $K=L^{O(1)}$ distinct pivot
vertices occur on minimum-$V$ paths from $x$ to $y$. The seed length,
the bit length of a path weight, and the workspace for evaluating a
vertex weight are all $O(L\ell)$.
\end{lemma}

\begin{proof}
We first construct a coloring for any fixed family of at most $N^2$
nonempty subsets of an $N$-element ground set. We want at most $K$
elements of the least occupied color in every set. At any stage retain
in each set the elements having its lexicographically least tuple of
colors so far. Suppose every retained set has size at most $2^a$.
Initially $a=L$. While $a>256\ell$, set
\[
           d=\lceil a/3\rceil+1,\qquad a'=\lceil 2a/3\rceil+2 .
\]
The next stage reduces the bound to $2^{a'}$. Choose a fresh
$\varepsilon$-biased string with $Nd$ binary coordinates using
Lemma~\ref{lem:bias} over $\mathbb F_2$. Give each vertex its own
block of $d$ bits. Its color is zero if all these bits are zero and
one otherwise. Take
\[
 t=2\lceil 64L/a\rceil,\qquad
 \varepsilon=2^{-10L-2t(a+1)}.
\]
This choice uses $O(L)$ seed bits since $t(a+1)=O(L)$.

Fix the choices of all preceding stages and consider one retained set
of size $2^{a'}<m\le2^a$. The expected number of zero-colored
vertices under full independence is $\mu=m2^{-d}$, with
$2^{a/3-1}\le\mu\le 2^{a'}/2$. Let $I_v$ indicate
that the first $d$ bits of vertex $v$ are zero, and let $X$ be their sum
over the retained set. Under fully independent bits these indicators
are independent with mean $p=2^{-d}$. In the expansion of
$\mathbb E[(X-\mu)^t]$, a tuple in which some vertex occurs exactly
once contributes zero. A remaining tuple has $j\le t/2$ distinct
vertices and contributes at most $(2p)^j$ in absolute value, since
$\mathbb E|I_v-p|^h\le 2p$ for all integers $h\ge 1$. There are at
most $m^j t^t$ tuples with $j$ distinct vertices. Hence the moment
under fully independent bits is at most
$2t^t(2\mu)^{t/2}$.

The Fourier coefficients of $I_v$ on the binary cube have absolute
sum one, by expanding the product of the $d$ indicators of a zero bit.
Thus the coefficients of $I_v-p$ have absolute sum at most two.
The coefficients of each product of $t$ centered indicators have
absolute sum at most $2^t$, including when vertices repeat. Under an
$\varepsilon$-biased distribution, its expectation therefore differs
from that under independence by at most $2^t\varepsilon$. Summing over
the $m^t$ tuples and applying Markov's inequality for even $t$ gives
\[
 \Pr[|X-\mu|\ge\mu]
 \le 2(2t^2/\mu)^{t/2}
      +\varepsilon(2m/\mu)^t
 \le 2^{-5L}
\]
for all sufficiently large $N$. Indeed, $a>256\ell$ and
$\log_2 t\le \ell+7$ give
$\log_2(\mu/(2t^2))\ge a/6$, while $t\ge128L/a$.
On the complementary event there are between one and $2^{a'}$
zero-colored vertices. All of them are retained. A set of size at most
$2^{a'}$ stays within the bound, irrespective of its colors. A union
bound over sets gives a seed that works for the entire stage.
Iterate with $a\gets a'$. There are $O(\ell)$ stages. The final bound
$K=2^{256\ell}$ is polynomial in $L$. Each stage uses $O(L)$ seed
bits and contributes a single binary color. The color tuple thus
encodes an integer $c(v)\in\{0,\ldots,C-1\}$ in lexicographic order
with $\log_2 C=O(\ell)$.

Apply the coloring to the nonempty feasible midpoint sets of the graph.
Put $B=4C+4$ and, at a positive layer $i$, set
\[
 V(v)=(1+c(v))B^{\nu_2(i)}.
\]
Take $V=0$ in layer zero. Fix an endpoint pair with pivot of valuation
$h$. For each $j\ge 1$ there are at most $2^j$ internal layers of
valuation $h-j$. Consequently their total contribution, summed over
all lower valuations along any path, is at most
\[
 C B^h\sum_{j\ge 1}(2/B)^j
 =\frac{2C}{B-2}B^h < B^h.
\]
An increase of one in pivot color costs $B^h$. Any minimum-$V$ path
must therefore use a feasible midpoint of the least occupied color.
There are at most $K$ such vertices. Finally $h\le L$, and a path has
at most $N$ vertices, giving bit length $O(L\ell)$. The colors and
the integer power can be evaluated within this space by regenerating
the relevant bits and using binary arithmetic.
\end{proof}
\section{Decomposing a frontier}

Fix successful initial weights $V$ from Lemma~\ref{lem:raw}. A \emph{window}
$I=(x,y)$ is a reachable ordered pair with at least one internal layer.
For $D\ge1$, give layer $t>0$ rank
\[
                   r_D(t)=\lfloor \nu_2(t)/D\rfloor .
\]
The \emph{frontier} $\fr_D(I)$ is the nonempty set of internal layer indices
of largest rank $r_D$ in the window. The $D$-skeleton of a path in $I$
lists the vertices it uses on these layers, in order. For a vector $W$
of weights, let $\calF_D(I;W)$ be the set of $D$-skeletons of $W$-optimal paths.
Write $f_D(I;W)=|\calF_D(I;W)|$. In particular, $f_1(I;V)\le K$.

Put $D_i=2^i$, for $0\le i\le s$, where $s$ is largest with
$D_s\le L/\ell$. Thus $s=O(\ell)$ and $D_s=\Theta(L/\ell)$.
We pass between successive frontiers as follows. Put $D=D_{i-1}$.
The binary valuations in a $2D$-frontier belong to one
band of $2D$ consecutive valuations. This band splits into a lower
half and an upper half of $D$ valuations each.
\begin{itemize}
\item First expose the $D$-skeleton of the entire window $I$.
If the greatest internal valuation belongs to the lower half, this
already specifies the $2D$-skeleton.
\item Otherwise the first exposure gives all upper-half vertices.
Between consecutive specified vertices, including the original endpoints,
expose the $D$-skeleton of each gap containing a lower-half frontier
layer. These exposures give precisely the remaining vertices of the
$2D$-skeleton.
\end{itemize}
Call each window on which this procedure asks for a $D$-skeleton a
\emph{site}. The site endpoints in the second step depend on the first
answer. On a fixed path the sites' $D$-frontiers partition $\fr_{2D}(I)$.
Site intervals can be nested, but their frontier layer sets are disjoint.

\section{Plans for naming the endpoints}

At refinement step $i$ the bound on the number of selected sites is
\[
                  m_i=\left\lceil\frac{L}{\ell D_i}\right\rceil .
\]
Fix an objective $W$ that lexicographically refines the raw weights $V$.
A history starts with one window and ends at any chosen scale. At step
$i$, choose a $W$-minimum path in each current parent window and select
at most $m_i$ of the sites in total from the decompositions along these
paths. The choices of sites in the same parent must all come from a
single path. Choices in different parents are independent, even if the
intervals of those parents overlap. For the next step, order the selected
sites first by the inherited order of their parents, then within each
parent by the order of the decomposition (the whole window first and
the gaps in layer order). Only the selected sites and their endpoints pass to the next step. At different refinement steps one may choose
new paths independently subject to the endpoint constraints already fixed.
Throughout a step the parent frontiers are pairwise disjoint. This holds
initially and follows at the next step from the partition within each
parent and the disjointness across parents.

Write $\calC_i(W)$ for the catalogue of all lists that occur at scale
$i$ in histories starting at scale $s$. In particular $\calC_s(W)$
contains all singleton windows. The lists realized by such histories are sufficient for the joint
induction. In its decision tree, one traversal corresponds to choosing
one optimal path for each parent window, and hence any selected subset
of its sites is compatible within each parent. Moreover, for any further
lexicographic refinement of $W$ the paths of that traversal are also
$W$-minimum. One appends the selected sites to a history realizing the
current parent list, using the prescribed order. This argument applies
both to the list of marked ambiguous sites and to the list of all
branching sites when its length is within the budget. The product of
frontier counts and the event of ties at all marked sites do not depend
on the ordering.

A structural plan records the initial endpoint pair, the stopping scale
and, at every step, the selected sites' parents and endpoint layers.
It does not specify any new vertex labels. In a parent band of $D_i$
consecutive binary valuations, the frontier contains at most $2^{D_i}-1$
layers. Indeed, all such layers are multiples of the lowest power of
two of the band between two consecutive multiples of the next band's
lowest power. There are at most $2^{D_i}+1$ possible endpoint layers
for a site, counting the parent endpoints. Their indices in layer order
use $O(D_i+1)$ bits each. The choices of layers recursively determine
all later layer positions, without reference to the intervening vertex
labels.

Here is an explicit encoding bound. At each step scan the current parents
in order and give the number of their selected children in unary, with a
terminator per parent. Then give the two endpoint-layer indices for each
child in its prescribed order. The unary strings cost the number of
parents plus the number of children. The total length of all the unary
strings in a history is $O(1+\sum_i m_i)$. Including the initial pair,
the stopping scale, and the indices, the length of a plan is
\[
 O\left(L+\sum_{i=1}^s m_i(D_i+1)\right)=O(L),
 \qquad
 \sum_{i=1}^s m_iD_i\le\frac{sL}{\ell}+\sum_{i=1}^sD_i=O(L).
\]
The empty choice at any refinement is allowed. Thus there are $2^{O(L)}$
structural plans. Equivalently one could spend $O(\ell)$ bits per child
to explicitly name its parent, since
$\sum_i m_i\ell\le L\sum_i 1/D_i+s\ell=O(L)$. In particular,
\begin{equation}\label{eq:pilot-budget}
                    \sum_{i=1}^s m_i(D_i+\ell)=O(L).
\end{equation}

For a fixed plan, order the queries for its vertex transcript
deterministically. For each parent reveal all requested endpoint vertices
jointly using one binary pivot recursion, taking the left side first.
A target that is already a parent endpoint requires no query.
Otherwise expose the pivot vertex, then recurse only into sides
containing an as yet unexposed target. Every target is on the parent's
$D_i$-frontier. Throughout its descent the pivot valuation strictly
decreases and stays at least that of the target until the target is
reached. Thus at most $D_i$ pivots are exposed for a given target,
all on the parent frontier. The total number of queries in the entire
history is bounded by
\[
                         2\sum_{i=1}^s m_iD_i=O(L).
\]
Within a parent, a common ancestor of different targets is exposed once.
Across parents in the same step, every exposed layer belongs to just one
parent frontier. Consequently a physical layer is queried at most once
in each refinement and at most $s\le\ell$ times in the whole history.
For fixed earlier answers the endpoints of the current binary window are
known. Every next answer extendible to a $W$-realizable history lies on
a $W$-minimum path for that binary window, by subpath replacement.
This last property gives at most $K$ initial raw choices per query, and
permits further thinning of the transcript tree by additional objectives.

\section{Thinning the trees of endpoint choices}

We first describe a perturbation $H$ of a fixed vector objective $W$.
It is a new, least significant lexicographic coordinate. Put the binary
valuations $0,\ldots,L$ in consecutive batches of at most $\Delta$
valuations, in increasing order. Let $p$ be prime and take $A$ to be
the least power of two at least $64p$. Independently in each batch $g$
choose $h_g\in\F_p^N$ using Lemma~\ref{lem:bias} with bias
$\varepsilon$. At a layer of valuation in batch $g$ and of position
$j$ within that batch (starting at zero), give vertex $v$ the weight
$H(v)=A^j[h_g(v)]_p$. At layer zero take weight zero.
Here $[a]_p$ denotes the representative in $\{0,\ldots,p-1\}$.

\begin{lemma}[Simultaneous midpoint ties]\label{lem:ties}
Fix pivot windows before drawing the batch seeds, one per distinct
physical pivot layer. In each window prescribe a center vertex and
one or more distinct alternatives at its pivot, all feasible on
$W$-optimal paths. If there are $R$ alternatives in total, the
probability that every prescribed center and alternative remains
feasible on $(W,H)$-optimal paths in its own window is at most
$p^{-R}+2(L+1)\varepsilon$, provided $(L+1)\varepsilon\le 1$.
\end{lemma}

\begin{proof}
For a window with top valuation at position $j$ in batch $g$, condition
on all earlier batch labels. For each prescribed pivot vertex $v$, let
$c_v$ be the minimum contribution of earlier batches among $W$-optimal
paths through $v$ in this window. There are no internal layers in later
batches. Within the current batch, the top layer contributes
$A^j[h_g(v)]_p$. The lower valuations contribute less than $A^j/8$
on every path, since the multiplicity $k$ valuations below the top
is at most $2^k$ and $p\sum_{k\ge1}(2/A)^k<1/8$. Consequently the
least $H$-cost in the class through $v$ is in
$[c_v+A^j[h_g(v)]_p,c_v+A^j[h_g(v)]_p+A^j/8)$.
If $u,v$ both occur on new minima, the integer
$[h_g(v)]_p-[h_g(u)]_p$ is within $1/8$ of $(c_u-c_v)/A^j$.
This gives at most one possible right-hand side fixed before the batch
is drawn. Reducing modulo $p$ produces an affine equation.
Within a window the differences from the center are linearly
independent, and different windows have disjoint pivot layers.
Thus if batch $g$ contains $R_g$ alternatives, the conditional
probability of its equations is at most $p^{-R_g}+\varepsilon$.
The equations of earlier batches are determined by the time batch $g$
is drawn. Multiplying the conditional bounds costs at most
$\prod_g(p^{-R_g}+\varepsilon)\le p^{-R}+(1+\varepsilon)^{L+1}-1
\le p^{-R}+2(L+1)\varepsilon$.
\end{proof}

Fix $Y=\lceil c_0L\rceil$ bounding the number of queries in any plan,
where $c_0$ is a sufficiently large absolute integer.
Under a fixed plan, let the branches at a query be the different next
vertex identities that are consistent with at least one complete
$W$-realizable history extending the current prefix. This gives a tree
with one leaf for every realizable vertex transcript. The window of the
current pivot query is determined by the plan and the earlier answers.
Each branch occurs on a $W$-optimal path in that window by subpath
replacement. For a leaf $P$, define
\[
        E_W(P)=\sum_{q\in P}\log_2 k_q,
\]
where $k_q$ is the number of branches at query $q$. Every leaf has
$E_V(P)\le Y\log_2 K$. More generally, if $E_W(P)\le Y\gamma$ for
every leaf of every plan, each plan has at most $2^{Y\gamma}$ vertex
transcripts. This follows by choosing a branch uniformly at each query.
A lexicographic refinement of $W$ only removes complete histories, so every
surviving leaf was also a leaf under $W$ with the same answers.

\begin{lemma}\label{lem:iter-pilot}
Under \eqref{eq:pilot-budget}, there is an explicitly enumerable family
of refinements $W$ of $V$ with total seed length and total path-weight
bit length $O(L\ell)$. At least one refinement satisfies
$|\calC_i(W)|\le 2^{O(L)}$ for every $i$. Its weights can be evaluated
in the same space given a seed.
\end{lemma}

\begin{proof}
Suppose $E_W(P)\le Y\gamma$ for every leaf, for an integer
$1\le\gamma\le O(\ell)$. We construct one further objective $H$ so
that $E_{(W,H)}(P)\le Y\gamma'$ with
\[
           \gamma'=\left\lceil\log_2(1+2(\gamma+1))\right\rceil.
\]
Choose a prime $p$ with $\log_2 p=\Theta(\ell)$, taking the lower
constant sufficiently large, and set
$\varepsilon=2^{-c_1Y(\gamma+1)}$ with $c_1$ sufficiently large.
Use the preceding batched perturbation, whose digit radix has
$\log_2 A=O(\ell)$.

Suppose a new leaf violates the proposed bound. Pad its sequence of
branch counts $k_q$ with ones to length $Y$. The arithmetic--geometric
mean inequality gives
\[
    \sum_q(k_q-1)\ge Y(2^{E_{(W,H)}(P)/Y}-1)>2Y(\gamma+1).
\]
At each query, connect the realized vertex to its $k_q-1$ alternatives.
These edges form a star of tied midpoint choices in the query's window. For each
physical layer retain a query with a largest such star. Since there
were at most $\ell$ queries per layer, the retained stars together have
at least $2Y(\gamma+1)/\ell$ alternatives. Keep exactly
$R=\lceil Y(\gamma+1)/\ell\rceil$ of these alternatives.

The resulting certificate is fixed by a structural plan, a $W$-realizable
transcript, a marker on selected queries, and the alternatives kept at
each marked query. The center of its star is the realized answer in
the transcript. Each alternative is named by its rank among the
at most $K$ raw $V$-optimal binary midpoints of that query's window.
Delimiters cost $O(Y+R)$ bits, so the number of certificates is at most
\[
  2^{O(L)+Y\gamma+O(Y+R)+R\lceil\log_2(K+1)\rceil}.
\]
This counting uses the older objective $W$, fixed before the random
labels of $H$ are chosen.

For a fixed certificate discard it if any named choice is infeasible on
$W$-optimal paths. For every remaining certificate, use the tie calculation of Lemma~\ref{lem:ties}
for every retained alternative and its star center. Conditional on earlier
batches, the baseline for each vertex is its minimum earlier-batch cost
over all $W$-optimal paths in its pivot window through that vertex.
No conditioning on the other answers of the transcript is needed for
this calculation. If the alleged leaf exists, each listed alternative
has some compatible completion in the new decision tree. Its subpath is
therefore $(W,H)$-optimal in this pivot window, as is the center's subpath.
Consequently the tie equations for all listed alternatives are necessary,
even when their completions differ. Within one star the coefficient rows
are the differences of distinct coordinate vectors and the center's
coordinate vector. These rows are linearly independent. Different retained
stars occupy different physical layers. All rows in a common batch are
therefore linearly independent. The proof of Lemma~\ref{lem:ties} bounds
the probability of all the required equations by
$p^{-R}+2(L+1)\varepsilon$.

Since $Y=\Theta(L)$, $R\ell\ge Y(\gamma+1)$, and
$\log_2(K+1)=O(\ell)$, choosing first the lower constant in
$\log_2 p$ and then $c_1$ makes the union bound less than one.
This proves the claimed reduction to $\gamma'$.

Start with $\gamma=\lceil\log_2(K+1)\rceil$ and repeat until $\gamma$
is bounded by an absolute constant. There are at most $L+1$ binary ranks.
Each batch seed in one iteration costs $O(L(\gamma+1))$ bits, and a radix
position costs $O(\ell)$ bits. Taking the batch length to be
$\lceil L\sqrt{(\gamma+1)/\ell}\rceil$, clipped to $[1,L+1]$, gives
total seed and weight bit length $O(L\sqrt{(\gamma+1)\ell})$ for that
iteration. More explicitly, storing the seeds costs
$O(L(\gamma+1)\lceil(L+1)/\Delta\rceil)$ bits, while the path weight
costs $O(L+\Delta\ell)$ bits. The clipping is harmless because
$\gamma=O(\ell)$. For all sufficiently large $\gamma$ the next value is
at most $\gamma/4$. The sum of these costs from an initial $\gamma=O(\ell)$
is $O(L\ell)$. Finally each structural plan has $2^{O(L)}$ transcripts,
and there are $2^{O(L)}$ plans. Every member of every compatible catalogue
is generated by at least one such history.
\end{proof}
\section{Packing all the main perturbations into one weight}

Fix a successful preliminary objective $V_*=(V,H_1,\ldots,H_r)$ given
by Lemma~\ref{lem:iter-pilot}. Write $\calC_i=\calC_i(V_*)$.
There are absolute constants $c_1,c_2$ such that
$|\calC_i|\le2^{c_1L}$ for every $i$ and $m_1\log_2 K\le c_2L$.
Take $c_2\ge1$. Since $V_*$ refines $V$, every window still has at
most $K$ possible optimal binary midpoint vertices.

Choose an absolute integer $C>c_1+2c_2+10$. Put
$d_i=D_{i-1}$, $b_i=C\ell D_i$, and choose the least prime $p_i>2^{b_i}$.
Set $\beta_i=b_i+1$, so that the integer representative of a label
in $\mathbb F_{p_i}$ is less than $2^{\beta_i}$. Choose an integer
$T=O(\ell)$ satisfying
\[
                     T d_i\ge\beta_i+2d_i+8\qquad(1\le i\le s).
\]
Set $E_0=0$ and $E_i-E_{i-1}=\beta_i+d_i+8$.
For a positive layer $t$ let
\[
    \sigma_i(t)=2^{\,T d_i\lfloor\nu_2(t)/d_i\rfloor-E_i}.
\]
Independently for each $i$ choose labels $g_i\in\mathbb F_{p_i}^N$
from Lemma~\ref{lem:bias} with
$\varepsilon=2^{-(c_1+2c_2+10)L}$.
For a vertex $v$ in positive layer $t$ define
\[
 U_i(v)=\sigma_i(t)[g_i(v)]_{p_i},\qquad
 P_i(v)=\sum_{a=1}^{i} U_a(v),\qquad
 M_i(v)=\sum_{a=i+1}^{s}\sigma_a(t)2^{\beta_a}.
\]
Take $P_0=0$ and set all these weights to zero on layer zero.
Here $P_i$ is a scalar sum. The values of $M_i$ are independent of all
labels and serve only in the proof.

For any window $I$ call a path \emph{admissible at step $i$} if it
minimizes $(V_*,P_i+e)$ lexicographically between its endpoints for
some vertex function $0\le e(v)\le M_i(v)$. Two admissible paths
may use different functions $e$. At step zero the admissible paths are
exactly the $V_*$-optimal paths, since $e=0$ is allowed. Admissibility
is nested: a witness $e$ at step $i$ gives the witness $U_i+e$ at step
$i-1$. It is also hereditary: replacing any subpath by a better path
with the same endpoints would improve the whole path under the same
objective and the same $e$. Write $f_D^{(i)}(I)$ for the number of
distinct $D$-skeletons of paths admissible at step $i$ in $I$.

\begin{lemma}\label{lem:packed-ties}
Fix the preceding labels and fix a list of $m$ windows with disjoint
$d_i$-frontiers. At each window prescribe two distinct $d_i$-skeletons
of paths admissible at step $i-1$. The probability that every prescribed
skeleton is still realized by a path admissible at step $i$ is at most
$p_i^{-m}+\varepsilon$.
\end{lemma}
\begin{proof}
Let $h$ be the largest binary valuation among the internal layers of
one window and put $u_a=d_a\lfloor h/d_a\rfloor$. Write
$\sigma=2^{T u_i-E_i}$. At a scale of width $d$, the number of internal
layers $j$ rank bands below the maximum band is at most $2^{(j+1)d}$
for every $j\ge0$: the window lies between consecutive multiples of the
first power of $2$ above the maximum band. Consequently the contribution
of $U_i$ below its frontier on any path is at most
\[
 \sigma 2^{\beta_i}\sum_{j\ge1}2^{(j+1)d_i-Td_i j}<\sigma/32.
\]
Since $d_i\mid d_a$ for $a>i$, we have $u_a\le u_i$. The contribution
from scale $a$ to the envelope $M_i$ on any path, including that
scale's frontier, is at most
\[
 2^{T u_a-E_a+\beta_a}\sum_{j\ge0}2^{(j+1)d_a-Td_a j}
 \le \sigma\,2^{\beta_a+d_a+1-(E_a-E_i)}
 \le \sigma\,2^{1-8(a-i)} .
\]
Summing over $a>i$ gives less than $\sigma/32$. Here and throughout the window comparison, path costs omit the fixed endpoints.

For a prescribed skeleton $S$ let $c_S$ be the minimum $P_{i-1}$-cost
over all \emph{$V_*$-optimal} paths with skeleton $S$ in the window.
This class is nonempty, and $c_S$ is fixed before $g_i$ is chosen.
Write $G_S=\sum_{v\in S}[g_i(v)]_{p_i}$. For any allowed $e$, the
least $(P_i+e)$-cost in this class belongs to
\[
               [c_S+\sigma G_S,\ c_S+\sigma G_S+\sigma/16].
\]
The lower bound follows by nonnegativity and the upper bound by using a
path attaining $c_S$. If $S$ is admissible, compare its class with the
class of the other prescribed skeleton $S'$ under a perturbation witnessing
$S$. This yields $c_S+\sigma G_S\le c_{S'}+\sigma G_{S'}+\sigma/16$.
Interchanging the skeletons, using a witness for $S'$, gives the reverse
inequality with the same tolerance. Thus the integer $G_S-G_{S'}$ is
within $1/16$ of the fixed number $(c_{S'}-c_S)/\sigma$. At most one
integer is possible. Reduction modulo $p_i$ gives a fixed affine equation
with nonzero coefficient row, or the event is impossible. Disjoint
frontier layers give these equations disjoint supports. The small-bias
bound proves the lemma.
\end{proof}

\paragraph{The product induction.}
We choose the labels sequentially so that, at every step $i$,
\[
          \prod_{I\in J} f_{D_i}^{(i)}(I)\le 2^{c_2 L}
                    \qquad (J\in\mathcal C_i).
\]
The initial claim follows from $m_1\log_2 K\le c_2 L$.
For any list of exactly $m_i$ sites in $\mathcal C_{i-1}$, the number
of ways to prescribe two old admissible skeletons per site is at most
$2^{2c_2 L}$ by induction. By Lemma~\ref{lem:packed-ties}, the probability
that some such list and prescription, with distinct alternatives at
every site, survive throughout is at most
\[
            2^{(c_1+2c_2)L}(p_i^{-m_i}+\varepsilon)<1,
\]
since $m_i b_i\ge C L$. Choose labels avoiding this event.

For a fixed parent list $J\in\mathcal C_i$ consider the decision tree
that exposes a tuple of admissible $D_i$-skeletons, with independent
choices of path in the different listed windows. Process sites in list
order and gaps in layer order, branching on each different answer
that extends to at least one complete admissible tuple. The tree has
one leaf for each tuple, even though a tuple can have many path
realizations. On any traversal the site frontiers are disjoint and all
site placements in a parent come from one $V_*$-optimal path. Thus any
selection of at most $m_i$ sites on the traversal, in the prescribed
canonical order, belongs to $\mathcal C_{i-1}$. If $m_i$ sites branch,
choose two alternatives at each. By heredity each alternative is
admissible in its own site, contradicting the event just excluded.
On every traversal the list of all branching sites is therefore in
$\mathcal C_{i-1}$. Choose uniformly among the available answers at
each decision node. Its branching factor is at most the old
$f_{d_i}^{(i-1)}(I)$ for that site $I$, by nesting and heredity.
Induction bounds the product of the branching factors on every traversal
by $2^{c_2 L}$. Every leaf has probability at least $2^{-c_2 L}$,
proving the product bound. This reasoning only uses answers with a
completion and does not assume that differently witnessed admissible
subpaths can be concatenated.

\paragraph{Final isolation.}
At $i=s$ the envelope is zero, so the admissible paths are exactly the
minimizers of $(V_*,P_s)$. Every singleton window belongs to
$\mathcal C_s$, hence has at most $2^{c_2L}$ optimal $D_s$-skeletons.
Choose a prime $p>2^{(2c_2+6)L}$ of $O(L)$ bits and one element
$z\in\mathbb F_p$. For two distinct such skeletons in one window,
the sums of the monomials $z^{\mathrm{id}(v)}$ differ as polynomials:
the positive vertex identifiers are distinct and each vertex occurs
at most once. Their difference has degree at most $N$. There are at
most $N^2 2^{2c_2 L}$ comparisons, so some $z$ separates all these
sums modulo $p$. Take the last lexicographic objective to give $v$ in
positive layer $t$ weight
\[
             (16Np)^{\lfloor\nu_2(t)/D_s\rfloor}
                     [z^{\mathrm{id}(v)}]_p .
\]
Set it to zero on layer zero. Within each window a unit difference in
the frontier sum dominates the entire contribution of lower ranks.
The winning skeleton is unique. Its gaps have smaller largest internal
rank, so induction shows that the entire winning path is unique, with
the empty-frontier case following from simplicity.

\paragraph{Resource bounds.}
Each main seed has $O(L+\beta_i)=O(L)$ bits, and there are $s=O(\ell)$
of them. Since $D_s=\Theta(L/\ell)$ and the widths double,
$E_s=O(\sum_i(\ell D_i+D_i+1))=O(L)$. Multiplying every $P_s(v)$ by
$2^{E_s}$ gives a nonnegative integer. Its path-cost bit length is
$O(TL+E_s+\max_i\beta_i+\log(sN))=O(L\ell)$.
The final objective has $O(1+L/D_s)=O(\ell)$ radix positions of
$O(L)$ bits. The bit lengths and seed costs of the preliminary
objectives also total $O(L\ell)$. Combining $V_*,2^{E_s}P_s$, and the
last objective using radices greater than their respective maximum
path costs yields integer weights whose path bit length is $O(L\ell)$.
Uniform upper bounds for the radices follow by multiplying the
maximum possible vertex cost in each coordinate by $N$ and adding one.
All constants are absolute. In particular $m_i$ is fixed before choosing
$c_1,c_2,C$, in that order of dependence. The primes can be found by
trial division. On input a seed and a vertex, field operations, binary
shifts, sums, and the mixed-radix encoding fit in the same space. The
catalogues and their decision trees enter only the existence proof.
\section{The unambiguous algorithm}
Enumerate all seed tuples in fixed order until the first min-isolating
tuple passes. This requires a check with one returning branch even for a
failed tuple. The weighted reachability checker of van Melkebeek and
Prakriya~\cite[Lemma 3.5]{vMP} uses space $O(\log N+\log(W+1))$ for an
upper bound $W$ on the path weights. For completeness, in the layered
case one can implement the checker as follows. Count all entered vertices
when weighing paths. At layer distance $j$ store the number of reachable
pairs and the sum of their minimum weights. For $j=0$ these numbers are
$N$ and zero. Suppose the census is correct and all the minima at distance
$j$ are unique. To look up any of these distances, scan all pairs at that
layer distance, guessing either to skip a pair or exhibit a path for it.
Abort unless the stored count and sum match. The count forces every
reachable pair to be represented. The sum then forces every exhibited path
to be minimum. Using a unique encoding of path guesses leaves exactly one
returning branch, which also knows the requested distance or its infinity.

For each distance-$(j+1)$ pair, scan the last edge and use these lookups
to find the best predecessor. Two predecessors tied for a finite minimum give a unique
return reporting a failed assignment. Otherwise accumulate the next
census and continue. This uses a constant number of counters of
$O(L+\log(W+1))$ bits. One can also record the designated pair's reachability,
handling same-layer and backward queries directly. Thus even the search
through failed assignments has a unique continuing branch. The isolation
argument ensures that the enumeration terminates.
This proves Theorem~\ref{thm:main}.

\paragraph{The cost of the initial palette.}
One limitation of the first weight $V$ can be seen on a three-layer
graph. In the middle layer take ordered vertices $v_1,\ldots,v_w$.
Include edges $x_i\to v_h$ for $i\le h$ and $v_h\to y_j$ for $h\le j$.
The nonempty feasible midpoint sets are precisely the intervals of
$[w]$. Suppose an ordered coloring uses $C$ colors and the least color
in every interval occurs at most $K$ times. Removing the globally least
color cuts the line into at most $K+1$ intervals, each using at most
$C-1$ colors. Induction gives $w\le(K+1)^C-1$. Thus for
$w=N^{\Theta(1)}$ and $K=L^{O(1)}$, the palette already needs
$\log_2 C\ge \ell-O(\log\ell)$ bits. The gadget can be placed with
middle layer $2^{\lfloor L/2\rfloor}$ by filling earlier layers with
isolated vertices. A common radix $B\ge C$ in the initial weight
$(1+c(v))B^{\nu_2(t)}$ then costs $\Omega(L\ell)$ bits. This gives a
specific obstacle to a further saving by shrinking the palette in
Lemma~\ref{lem:raw}.

\begin{thebibliography}{9}
\bibitem{vMP} Dieter van Melkebeek and Gautam Prakriya.
\emph{Derandomizing isolation in space-bounded settings}.
SIAM Journal on Computing 48(3), 979--1021, 2019.
\href{https://doi.org/10.1137/17M1130538}{doi:10.1137/17M1130538}.
\bibitem{Hoza} William M. Hoza.
\emph{Deterministic, oblivious isolation for space-bounded computation requires large weights}.
ECCC TR26-124, 2026. \url{https://eccc.weizmann.ac.il/report/2026/124/}
\bibitem{RA} Klaus Reinhardt and Eric Allender.
\emph{Making nondeterminism unambiguous}.
SIAM Journal on Computing 29(4), 1118--1131, 2000.
\href{https://doi.org/10.1137/S0097539798339041}{doi:10.1137/S0097539798339041}.
\end{thebibliography}
\end{document}
